\documentclass[11pt]{article}

\usepackage{sectsty}
\usepackage{graphicx}
\usepackage{xcolor}

\usepackage[round]{natbib}

\usepackage{framed}

\usepackage[normalem]{ulem}
\usepackage[hidelinks]{hyperref}
\usepackage{hyperref}
\hypersetup{
  colorlinks=true,
  citecolor=black,
  citecolor=violet,
  linkcolor=black,    % internal refs like \ref, \pageref — optional
  urlcolor=blue        % keeps \href links visibly blue/clickable
}
\let\oldhref\href
\renewcommand{\href}[2]{\oldhref{#1}{\uline{#2}}}
%\usepackage{url}
% Margins
\topmargin=-0.45in
\evensidemargin=0in
\oddsidemargin=0in
\textwidth=6.5in
\textheight=9.0in
\headsep=0.25in

\usepackage{parskip}

\title{\vspace{-0.5in}Reply to Reviewers\vspace{-0.5in}}
\author{}
%\vspace{-1in}
\date{\today}


\begin{document}
\maketitle	

{\color{violet}
\begin{leftbar}
In the revised manuscript, changes made in response to the first reviewer's comments are labeled with r1; changes made in response to the second reviewer's comments are labeled with r2. Typographical errors corrected by the authors are labeled with a.
\end{leftbar}
}

\section{Reviewer 1}

{\color{violet}
\begin{leftbar}
We thank the reviewer for their careful reading and suggestions. In the following, we provide details of the changes made.
\end{leftbar}
}

Reviewer \#1 (Formal Review for Authors (shown to authors)):
\\
\\
This is the second round of review for this manuscript. Examining carefully the authors' replies and the modifications made for the text, the reviewer confirms that the authors have fully addressed the comments given in the previous review. The reviewer is happy to recommend the manuscript for publication now.

One minor editorial concern is that in the Figure 3 that the authors uploaded, the right edges of the legends for the right-column plots seem to be cut off somehow. They should properly be fixed in the production process of a camera-ready version of the paper.
\pagebreak

\section{Reviewer 2}

\subsection{Overall remarks and recommendations}

The authors have done a great job of revising the manuscript. I would like to thank them for the way that they have addressed my many comments in a positive way. I must also apologise for taking so long to complete my review; it’s been a crazy summer.

Looking in more detail at the revised manuscript, I am pleased to note that the revised abstract, and the addition of a plain language summary, both now highlight that the manuscript aims to promote the development of up-to-date and robust standards for space physics coordinate systems. In parallel, the use of SPICE is now clearly presented as a widely used tool that can provide a platform for implementing those
standards. I was also very pleased to see how the authors have updated the manuscript to include recommendations that the proposed standards should be developed internationally and present an excellent initial list of international groups that could contribute to this development.

I also thank the authors for responding to my comment on the need for more historical background on how space physics coordinate system developed from astronomical systems, particularly right ascension and declination. They have made some updates to address this in the current manuscript, notably at new lines 168 to 169, where the key association between GEI and right ascension/declination is noted. This will help
readers, like myself, who have come into space physics with a background in observational astronomy. The authors have also suggested that a wider review of this historical background should form part of the proposed development of standards for space physics coordinate systems. I am very happy to concur with their suggestion.

I also thank the authors for providing a clear justification of why the detailed analysis in section 4 (Comparisons) should be in the main paper, rather than Supplementary Information, as I had originally suggested. I note that this is reinforced by text changes in the Introduction that add weight to the role of
section 4.

In summary, I consider that this manuscript now provides a strong and clear justification of the need to develop new standards for space physics coordinate transformations. I anticipate that the present manuscript, when published, will be a key starting point for future work on standards for these coordinate systems. Thus I am keen to see this manuscript published. However, there are four minor points that I would ask the authors to consider before publication. These are listed below. One (item 3) is a small change to reduce a risk of misinterpretation. The others are small changes that will reinforce the value of the manuscript as a starting point for future work.

Minor comments

1. Line 159-160. The revised manuscript stimulated me to look at Simon Newcomb’s 1898 report and at the 1961 Explanatory Supplement from HMNAO. As best I can see the GMST formula used in Russell (1971) is explicitly stated only in the 1961 Explanatory Supplement, but with an explanation on pages 74-75 of how it is derived from other formulae given on page 9 of Newcomb’s report. So I recommend that the authors
include those page numbers in the reference entries for HMNAO (1961) and Newcomb (1898). I think this will reinforce the future value of this paper as a starting point for future work in this area. (As an aside I note that the formulae in both papers continue to work well, e.g. giving GMST to one second accuracy in August 2026.)

2. Lines 574-581. I thank the authors for including this discussion of time standards, and particularly the current close alignment of UT1 and UTC, which has been in place since 1972. However, I strongly encourage the authors to note that this will soon end, certainly by 2035, maybe earlier (Gibney, 2022). This is an important issue for the future work on, and use of, space physics coordinate systems. For that reason I encourage the authors to briefly note this issue in the current manuscript.

3. Lines 692-694. The sentence beginning “In general, the mission has access to positions …” could be misinterpreted as implying that the Cluster spacecraft carried on-board GNSS. But they were designed more than 30 years ago, before the general adoption of on-board GNSS as a satellite tracking technology in the late 1990s and early 2000s. So to avoid this misinterpretation, I suggest to start the sentence “In general, modern missions have access to positions …”.

4. Line 679. I was also surprised that the authors omitted to include ranging as a key tracking technique used by satellite operators. I should perhaps have called it “active ranging”, i.e. where the return signal is generated by a transponder on the spacecraft, thus giving a strong return that can be measured with high accuracy. This technique is, of course, limited to the case where the satellite is designed to cooperate with the tracking system on the ground. But that is the case for space physics missions. I leave
it to the authors to decide whether to incorporate this minor point in the present manuscript. I raise it here to encourage them to take account of active ranging, whether based on radar or laser tracking, in their future work.

{\raggedright
\bibliographystyle{apalike}
\bibliography{main}
}

\end{document}
